\section{Dialoge}

\subsection{Die Umgebung \texttt{speech}}

Die Umgebung \verb|speech| dient zur Darstellung von gesprochener Sprache.

Der Inhalt wird

\begin{itemize}
\item kursiv gesetzt,
\item in normaler Schrift dargestellt,
\item und übernimmt Zeilenumbrüche exakt aus dem Quelltext.
\end{itemize}

Dies eignet sich besonders zum Vorlesen von Dialogen oder Monologen.

\subsubsection*{Syntax}

\begin{Verbatim}
\begin{speech}
...
\end{speech}
\end{Verbatim}

\subsubsection*{Beispiel}

\paragraph{Quellcode}

\begin{Verbatim}
\begin{speech}
Willkommen in meinem Gasthaus.

Setzt euch.
Das Feuer brennt bereits.
\end{speech}
\end{Verbatim}

\paragraph{Ergebnis}

\begin{speech}
Willkommen in meinem Gasthaus.

Setzt euch.
Das Feuer brennt bereits.
\end{speech}
